% !TeX root = ../机器学习笔记.tex
% 由主文档通过 \input 载入；请编译主文档。
\setcounter{chapter}{22} % 保留原书章号。
\chapter{半监督、主动学习与迁移学习}
\label{chap:ml-transfer}

标签稀缺和分布变化使学习问题更接近实际，也引入新的假设。额外数据只有通过合适的结构与评价条件才能提供帮助，数量本身不是保证。

\section{有限标签与半监督学习}

半监督学习同时使用标注集 $\mathcal D_L$ 和未标注集 $\mathcal D_U$。常见假设是附近样本标签相似、类别边界穿过低密度区域，或同一数据簇主要属于同一类别。在这些假设成立时，未标注数据帮助识别输入分布结构，从而限制分类边界。

如果标签恰好沿高密度簇内部变化，按簇传播标签会强化错误；如果未标注数据含很多不属于目标任务的类别，更多数据还可能损害训练。因而应先比较只有标注数据的基线，再评估额外结构是否符合任务，而不能把未标注数据视为无条件有益。

\section{伪标签与一致性训练}

伪标签（pseudo-labeling）以当前或教师模型为未标注输入产生 $\widehat y(x)=\arg\max_cp_c(x)$，只对最大概率超过阈值的样本加入监督式损失，标签在该更新中视为固定。高置信错误会自我强化，因此阈值、类别分布和教师更新方式需要独立验证。

一致性正则化（consistency regularization）要求同一输入的两种允许变换 $T_1,T_2$ 产生接近的预测，可使用平方差或 KL 损失。其前提是变换保留目标语义；若图像旋转改变类别、文本删词改变否定意义，强迫一致就可能错。未标注项通常另乘权重 $\lambda_U$，以控制错误信号对标注任务的影响。

\section{主动学习与标注预算}

主动学习（active learning）从未标注池中选择要付出代价获取标签的样本。分类不确定性可用预测熵 $-\sum_cp_c(x)\log p_c(x)$ 或前两类概率差衡量；代表性方法则避免只查询孤立异常点。查询后获得真实标签、重新训练，再重复直到标注预算耗尽。

模型不确定不一定意味着标签有价值：异常输入可能一直难以学习，过度自信区域也可能隐藏系统性错误。有效评价以相同初始标注集、同一标注预算和独立测试集比较主动查询与随机查询，并考虑每次重训与查询的计算成本。

\section{迁移学习与领域适应}

迁移学习（transfer learning）利用源任务知识改善目标任务，常见方式是冻结预训练表示、只训练预测头，或进一步微调部分乃至全部参数。源任务相似性、目标数据量和模型容量决定策略取舍；源表示不含目标所需信息时，迁移可产生负收益，即负迁移。

领域适应（domain adaptation）着重处理源与目标分布差异。必须说明目标标签、未标注目标输入及源数据是否可用；使用目标测试输入进行传导适应与完全未见目标样本的归纳测试是不同设定。预训练语料与测试样本重叠，也会改变泛化结论的解释。

\section{分布变化与重要性加权}

协变量偏移（covariate shift）假设 $p_t(y\mid x)=p_s(y\mid x)$，而输入边缘分布改变。若 $p_t(x)>0$ 的地方都有 $p_s(x)>0$，定义 $w(x)=p_t(x)/p_s(x)$，则
\begin{equation}
 \mathcal R_t(f)=\int\ell(f(x),y)p_t(x)p_s(y\mid x)\,dx\,dy
 =\mathbb E_s[w(X)\ell(f(X),Y)].
 \label{eq:ml-covariate-shift}
\end{equation}
这给出源样本加权评价目标风险的依据，但要求损失可积。目标支持超出源支持时，源数据根本没有对应观测，权重修正无法凭空补足。

类别先验变化假设类条件分布保持不变，条件关系变化则直接改变 $p(y\mid x)$。不同变化不能使用同一修正。估计密度比可能很难，极端权重还会放大方差，截断权重则引入偏差；应明确相应取舍。

\section{例题、练习与学习检查}

\begin{example}[离散协变量偏移]
输入取 $a,b$，源概率为 $(0.8,0.2)$，目标概率为 $(0.4,0.6)$，两处条件损失期望分别为 $0.1,0.5$。
\end{example}
\begin{solve}
源风险为 $0.18$，目标风险为 $0.34$。权重分别为 $0.5,3$，加权源风险为 $0.8(0.5)(0.1)+0.2(3)(0.5)=0.34$。如果 $b$ 在源分布中概率为零，就不能进行这种重加权恢复。
\end{solve}

习题：构造未标注簇结构误导分类的例子；设计按标签预算比较的主动学习曲线；区分协变量偏移、标签先验偏移和概念变化；解释为什么低熵伪标签仍可能系统性错误。





